% !TEX root = main.tex

\section{Introducción a la capa de enlace}

La presente capa permite a las capas superiores acceder a los medios, acepta paquetes de la capa 3 y los empaqueta en \textit{tramas}. Además de preparar los datos de red para la red fisica, este controla la forma en que los datos se ingresan y se reciben en los medios.

\subsection{Nodos}
Los nodos crean y reenvían tramas.
\subsection{Subcapas de enlace de datos}
\begin{enumerate}
\item Control de enlace Lógico (LLC): esta subcapa superior se comunica con la capa de red
\item Control de acceso al medio (MAC): define los procesos de acceso al medio que realiza el hardware.
\end{enumerate}

\subsection{Diversos canales de acceso múltiple}
Hay diversos canales como \textbf{Shared wire: cable access network, Shared wireless: WiFi, Satellite y Cocktail party}

El \textbf{Control de acceso a los medios} consiste en que los protocolos en la capa de enlace de datos definen las reglas de acceso a los diferentes medios.

Estas técnicas de control de acceso a lso medios definen si los nodos comparten los medios y de que manera lo hacen.

\subsection{Topología fisica}
Se refiere a las conexiones fisicas e identifica como se interconectan los nodos de la infraestructura. e.g (estrella, punto a punto, arbol), etc.

\begin{enumerate}
\item La \textbf{topología lógica}: se refiere a la forma en que una red transfiere tramas de un nodo al siguiente. Esta disposición consta de conexiones virtuales entre los nodos de una red. Los protocolos de capa de enlace de datos definen estas rutas de señales logicas.



\end{enumerate}

En la vida real, hay dos servicios que ofrece la capa de enlace en una red local:

\begin{enumerate}
\item 
\end{enumerate}
